\documentclass[11pt]{article}
\usepackage{amsmath,amssymb,amsthm}
\usepackage[margin=2.5cm]{geometry}

\newtheorem{theorem}{Theorem}
\newtheorem{lemma}[theorem]{Lemma}
\newtheorem{corollary}[theorem]{Corollary}
\theoremstyle{remark}
\newtheorem{remark}[theorem]{Remark}
\theoremstyle{definition}
\newtheorem*{statement}{Official statement}

\title{Problem 4, Cell 1: Sun's conjecture for $k=3$}
\author{}
\date{}

\begin{document}
\maketitle

\begin{statement}
For integers $a$ and $m\ge1$ write $a\ (\mathrm{mod}\ m)=\{a+mt:t\in\mathbb{Z}\}$. A finite family
of classes $a_1\ (\mathrm{mod}\ m_1),\dots,a_k\ (\mathrm{mod}\ m_k)$ is pairwise disjoint if
$a_i\ (\mathrm{mod}\ m_i)\cap a_j\ (\mathrm{mod}\ m_j)=\emptyset$ whenever $i<j$. Nothing is assumed
about the moduli beyond $m_i\ge1$: they may repeat, and the residues $a_i$ may repeat as well. By the
Chinese remainder theorem,
\[
  a_i\ (\mathrm{mod}\ m_i)\cap a_j\ (\mathrm{mod}\ m_j)\neq\emptyset\iff\gcd(m_i,m_j)\mid a_i-a_j.\tag{$*$}
\]
\emph{Question.} Let $k\ge2$ and let $a_1\ (\mathrm{mod}\ m_1),\dots,a_k\ (\mathrm{mod}\ m_k)$ be
pairwise disjoint. Must there be a pair $i<j$ with $\gcd(m_i,m_j)\ge k$?

The first size that the two trivial observations above do not settle. Prove the statement for
$k=3$: any three pairwise disjoint classes have $\gcd(m_i,m_j)\ge3$ for some $i<j$.
\end{statement}

\noindent Theorem~\ref{thm:main} below is exactly this statement. The notation $a\bmod m$ below is the class
$a\ (\mathrm{mod}\ m)$ of the statement. Lemma~\ref{lem:crt} is $(*)$, proved for completeness.

\paragraph{Known result.} The question is Sun's disjoint congruence classes conjecture~\cite{Sun06}.
O'Bryant~\cite{OBr07} proved it for all $k\le20$, so this case is known. The proof below is an
independent proof and is self-contained.

For an integer $m\ge1$ and $a\in\mathbb{Z}$ write $a\bmod m=\{x\in\mathbb{Z}: x\equiv a\pmod m\}$.
A family of residue classes is \emph{disjoint} if its members are pairwise disjoint.

\begin{theorem}\label{thm:main}
Let $a_1\bmod m_1,\ a_2\bmod m_2,\ a_3\bmod m_3$ be three pairwise disjoint residue classes.
Then $\gcd(m_i,m_j)\ge3$ for some $i\ne j$.
\end{theorem}

\section{When do two classes meet?}

\begin{lemma}\label{lem:crt}
Let $m,n\ge1$, $a,b\in\mathbb{Z}$ and $g=\gcd(m,n)$. Then
$(a\bmod m)\cap(b\bmod n)\ne\emptyset$ if and only if $g\mid a-b$.
\end{lemma}

\begin{proof}
($\Rightarrow$) If $x\equiv a\pmod m$ and $x\equiv b\pmod n$, then $g\mid m$ and $g\mid n$ give
$a\equiv x\equiv b\pmod g$.

($\Leftarrow$) By B\'ezout there are $u,v\in\mathbb{Z}$ with $um+vn=g$. Write $a-b=gt$ with
$t\in\mathbb{Z}$ and set $x=a-umt$. Clearly $x\equiv a\pmod m$. Also
$x=a-(g-vn)t=a-gt+vnt=b+vnt\equiv b\pmod n$.
\end{proof}

\begin{corollary}\label{cor:small}
Let $a\bmod m$ and $b\bmod n$ be disjoint and $g=\gcd(m,n)$. Then $g\ge2$ and $a\not\equiv b\pmod g$.
In particular, if $g=2$, then $a\not\equiv b\pmod 2$.
\end{corollary}

\begin{proof}
By Lemma~\ref{lem:crt}, disjointness means $g\nmid a-b$. Since $1$ divides every integer, $g\ne1$.
\end{proof}

\section{Proof of Theorem~\ref{thm:main}}

Suppose, for contradiction, that $\gcd(m_i,m_j)\le2$ for all $i\ne j$. By
Corollary~\ref{cor:small} each of these gcds is at least $2$, so $\gcd(m_i,m_j)=2$ for all three
pairs. By Corollary~\ref{cor:small} again, $a_i\not\equiv a_j\pmod2$ for all $i\ne j$. In other
words, $a_1,a_2,a_3$ lie in three distinct residue classes modulo $2$. There are only two such
classes, so this is impossible by the pigeonhole principle. Hence some pair satisfies
$\gcd(m_i,m_j)\ge3$. \qed

\begin{remark}
The bound is sharp: $0\bmod3$, $1\bmod3$, $2\bmod3$ are disjoint and every pairwise gcd equals $3$.
The case of a modulus $m_i=1$ needs no separate treatment: then $a_i\bmod m_i=\mathbb{Z}$, which meets
every class, and the corresponding gcds equal $1$, which Corollary~\ref{cor:small} excludes.
\end{remark}

\begin{thebibliography}{9}
\bibitem{OBr07} K.~O'Bryant, \emph{On Z.-W.\ Sun's disjoint congruence classes conjecture}, in
  \emph{Combinatorial Number Theory}, de Gruyter, Berlin, 2007, pp.~403--411; arXiv:math/0604347.
\bibitem{Sun06} Z.-W.~Sun, \emph{Finite covers of groups by cosets or subgroups}, Internat.\ J.\ Math.\
  \textbf{17} (2006), 1047--1064; arXiv:math/0501451.
\end{thebibliography}

\end{document}
